% Part: first-order-logic
% Chapter: natural-deduction
% Section: rules-and-proofs

\documentclass[../../../include/open-logic-section]{subfiles}

\begin{document}

\iftag{FOL}
      {\olfileid{fol}{ntd}{rul}}
      {\olfileid{pl}{ntd}{rul}}

\olsection{বিধি ও \usetoken{P}{derivation}}

\begin{explain}
স্বাভাবিক নিষ্পাদন-পদ্ধতিগুলি গাণিতিক প্রমাণে ব্যবহৃত অনানুষ্ঠানিক
যুক্তিবিচারের খুব কাছাকাছি গড়ে তোলা হয় (এই অর্থেই পদ্ধতিটি কিছুটা
``স্বাভাবিক'')। স্বাভাবিক নিষ্পাদনের প্রমাণ অনুমিতি দিয়ে শুরু হয়।
তারপর অনুমান-বিধি প্রয়োগ করা হয়। \Intro{\lnot}, \Intro{\lif},
\iftag{FOL}{\Elim{\lor} ও \Elim{\lexists}}{ এবং \Elim{\lor}}
অনুমান-বিধি অনুমিতিগুলিকে ``!!{discharged}'' করে; স্পষ্টতার জন্য
!!{discharged} অনুমিতির চিহ্নটি অনুমানের পাশে বসানো হয়।
\end{explain}

\begin{defn}[অনুমিতি]
কোনো শাখার সর্বোচ্চ স্থানে থাকা যেকোনো !!{sentence} হলো একটি
\emph{অনুমিতি}।
\end{defn}

স্বাভাবিক নিষ্পাদনের !!^{derivation}s হলো !!{sentence}s-এর বিশেষ ধরনের
বৃক্ষ। এর সর্বোচ্চ !!{sentence}s অনুমিতি; আর কোনো !!a{sentence}-এর
ওপরে এক, দুই বা তিনটি অন্য বাক্য থাকলে সেটিকে কোনো অনুমান-বিধি থেকে
সঠিকভাবে অনুসৃত হতে হবে। অনুমানটির ওপরে থাকা !!{sentence}s-কে তার
\emph{পূর্বধারণা} এবং নিচের !!{sentence}-টিকে তার \emph{সিদ্ধান্ত}
বলা হয়। প্রত্যেক !!{operator}-এর জন্য বিধিগুলি জোড়ায় আসে—একটি
প্রবর্তন-বিধি ও একটি অপসারণ-বিধি। এগুলি সিদ্ধান্তে কোনো
!!a{operator} প্রবর্তন করে অথবা বিধির কোনো পূর্বধারণা থেকে
!!a{operator} অপসারণ করে। কয়েকটি বিধি বিশেষ ধরনের অনুমিতিকে
\emph{!!{discharged}} করার অনুমতি দেয়। কোন অনুমান কোন অনুমিতিকে
!!{discharged} করে তা বোঝাতে অনুমিতি ও অনুমান—দুটিকেই চিহ্ন দিই।
অনুমিতিটিকে ``$\Discharge{!A}{n}$'' লিখে এটি দেখানো হয়।
% BN-SRC-033 TARGET-CORRECTION-BEGIN
\par\smallskip\noindent\textnormal{\small\textbf{উৎস-সংশোধন (BN-SRC-033):} হিমায়িত ইংরেজি গদ্যে বৃক্ষের কোনো বাক্যের ওপরে থাকা এক, দুই বা তিনটি নোডকে “other sequents” বলা হয়েছে। এই অধ্যায়ের সংজ্ঞা ও প্রদর্শিত বৃক্ষগুলিতে নোডগুলি বাক্য, সিকোয়েন্ট নয়; তাই বাংলা গদ্যে “অন্য বাক্য” ব্যবহার করা হয়েছে।} % BN-SRC-033 TARGET-CORRECTION-END

% Only include this sentence is on of \land, lor, lif or lnot is defined.

\iftag{notprvNot,notprvAnd,notprvOr,notprvIf}{}%
	{$\land$, $\lor$, $\lif$, $\lnot$ ও $\lfalse$-এর কোনোটি সংজ্ঞায়িত হলেও, প্রচলিতভাবে সব !!{operator}s-এর বিধিই বিবেচনা করা হয়।}



\end{document}
